% !TeX root = ../main.tex
\chapter{JADWAL PENELITIAN}

\section{Rencana Kerja}

Jadwal berikut merupakan rencana tentatif enam bulan sejak penelitian disetujui. Nomor bulan digunakan agar jadwal tidak mengunci penelitian pada tanggal yang belum disepakati. Setiap perubahan setelah bimbingan akan dicatat pada versi berikutnya.

\begin{table}[htbp]
\centering
\renewcommand{\arraystretch}{1.25}
\setlength{\tabcolsep}{4pt}
\newcommand{\taskbar}{\cellcolor{blue!30}}
\begin{tabular}{|p{5.2cm}|c|c|c|c|c|c|}
\hline
\textbf{Kegiatan} & \textbf{B1} & \textbf{B2} & \textbf{B3} & \textbf{B4} & \textbf{B5} & \textbf{B6} \\
\hline
Tinjauan pustaka dan pematangan rumusan masalah & \taskbar & \taskbar & & & & \\
\hline
Audit sumber data dan penyusunan pedoman anotasi & \taskbar & \taskbar & & & & \\
\hline
Pilot annotation dan pengukuran IAA & & \taskbar & \taskbar & & & \\
\hline
Implementasi pipeline dan tiga konfigurasi eksperimen & & \taskbar & \taskbar & \taskbar & & \\
\hline
Pengujian pada validation set dan pembekuan protokol & & & \taskbar & \taskbar & & \\
\hline
Eksperimen test set dan pencatatan log & & & & \taskbar & \taskbar & \\
\hline
Analisis metrik, ablation, dan failure mode & & & & & \taskbar & \taskbar \\
\hline
Penulisan laporan dan konsultasi hasil & & & & & \taskbar & \taskbar \\
\hline
\end{tabular}
\end{table}

\noindent\textit{Keterangan:} B1--B6 adalah bulan relatif sejak persetujuan rencana penelitian. Blok biru menunjukkan periode aktivitas utama; beberapa aktivitas dapat berjalan paralel apabila dependensinya telah terpenuhi.
